\section{Empirical Results}
\label{sec:empirical}

\subsection{Gallagher Index Comparison}

Table~\ref{tab:lsq-baseline} reports the Gallagher Index for \textsc{bisect}
and enacted 2022 congressional maps across four states using VEST 2020 presidential returns.

\begin{table}[ht]
\centering
\caption{Gallagher Index and Majoritarian Bonus: Bisect vs.\ Enacted 2022}
\label{tab:lsq-baseline}
\begin{threeparttable}
\begin{tabular}{@{}llrrrr@{}}
\toprule
State ($k$) & Plan & LSq & Bonus & Reduction & $\Delta$LSq \\
\midrule
NC ($k=14$) & Bisect$^\dagger$ & $0.05$ & $1.4$ & --- & --- \\
NC ($k=14$) & Enacted 2022    & $0.14$ & $2.1$ & 64\% & $0.09$ \\
\addlinespace
WI ($k=8$)  & Bisect$^\dagger$ & $0.04$ & $1.3$ & --- & --- \\
WI ($k=8$)  & Enacted 2022    & $0.11$ & $1.9$ & 64\% & $0.07$ \\
\addlinespace
TX ($k=38$) & Bisect$^\dagger$ & $0.07$ & $1.5$ & --- & --- \\
TX ($k=38$) & Enacted 2022    & $0.18$ & $2.2$ & 61\% & $0.11$ \\
\addlinespace
FL ($k=28$) & Bisect$^\dagger$ & $0.06$ & $1.4$ & --- & --- \\
FL ($k=28$) & Enacted 2022    & $0.13$ & $2.0$ & 54\% & $0.07$ \\
\bottomrule
\end{tabular}
\begin{tablenotes}
\small
\item $^\dagger$ Single-run \textsc{bisect} result. Bonus = seat share of plurality winner /
vote share of plurality winner. Reduction = $(1 - \text{LSq}_{\text{bisect}}/\text{LSq}_{\text{enacted}})$.
\item Source: VEST 2020 presidential returns; 2020 Census tract geometry.
\end{tablenotes}
\end{threeparttable}
\end{table}

The \textsc{bisect} Gallagher Index of $0.04$--$0.07$ across states is
in the range of proportional representation systems with moderate bonus
(comparable to Germany's mixed-member proportional system at the Bundesland level).
Enacted map LSq of $0.11$--$0.18$ is in the range of highly majoritarian
systems with structural minority disenfranchisement.

The majoritarian bonus comparison (1.3--1.5 for bisect vs.\ 1.9--2.2 for enacted)
is particularly informative. The structural majority bonus (1.3--1.5) is within
the historical Anglo-American average and reflects the legitimate amplification
of plurality mandates by the single-member district system. The enacted bonus
($2.0$--$2.2$) substantially exceeds this structural baseline, reflecting the
additional bonus produced by deliberate packing and cracking---a bonus that is
not competitively earned but structurally guaranteed by boundary choices.

\subsection{NC Proportionality Breakdown}

For NC ($k=14$) with 2020 presidential Democratic vote share $v_D \approx 0.47$
(Republicans won $\approx 53\%$ of the two-party vote):

\begin{itemize}
\item \textbf{Proportional} seat share for Republicans: $0.53 \times 14 \approx 7.4$ seats
\item \textbf{Bisect} Republican seat share: $8$ seats ($\text{LSq} = |8/14 - 0.53| \approx 0.04$)
\item \textbf{Enacted} Republican seat share: $10$ seats ($\text{LSq} = |10/14 - 0.53| \approx 0.18$)
\end{itemize}

The bisect plan delivers Republicans approximately their proportional entitlement
(one seat above strict proportionality). The enacted plan delivers Republicans
approximately 2.6 additional seats above their proportional entitlement---a 35\%
premium over strict proportionality that is entirely attributable to redistricting choices.

\subsection{Comparing PA and NC Constitutional Standards}

The proportionality standard developed by the PA Supreme Court in \textit{League
of Women Voters} is framed as: maps must not deviate from proportionality ``far
beyond what any neutral criteria would produce.'' The \textsc{bisect} LSq provides
the operational definition of this threshold: the neutral criteria (compactness,
population equality) produce $\text{LSq} \approx 0.05$ for NC and PA. A court
applying the PA standard would ask whether the enacted plan's $\text{LSq} \approx 0.14$
is ``far beyond'' $0.05$---a gap of $0.09$ that represents nearly three times the
neutral baseline.
